\documentclass[11pt,oneside]{article}

\pdfminorversion=7
\usepackage{QuizStyle}

\hypersetup{
    pdftitle={20260916 Quiz for MATH200 Differential Equations},
    pdfauthor={Jungwoo Kim},
    pdfsubject={Quiz for MATH200 Differential Equations},
    pdfcreator={LaTeX}
}

\begin{document}

\quiztitle{20260916 Quiz}

% Boyce Section 2.1, Problem 3(c), treated here as the complete quiz problem.
\begin{problem}{2.1.3}
Find the general solution of the differential equation
\[
    y'+y=te^{-t}+2,
\]
and use it to determine how solutions behave as $t\to\infty$.
\end{problem}

\begin{solution}
The integrating factor is
\[
    \mu(t)=\exp\!\left(\int 1\,\dd t\right)=e^t.
\]
Multiplying the equation by $e^t$ gives
\[
    (e^t y)'=t+2e^t.
\]
Hence
\[
    e^t y=\frac{t^2}{2}+2e^t+c,
\]
so the general solution is
\[
    \boxed{y=2+\left(\frac{t^2}{2}+c\right)e^{-t}}.
\]
Since $t^2e^{-t}\to0$ and $e^{-t}\to0$ as $t\to\infty$, every solution satisfies
\[
    \boxed{y(t)\to2\quad\text{as }t\to\infty}.
\]
\end{solution}

\begin{problem}{2.1.12}
Find the solution of the initial value problem
\[
    ty'+2y=2\sin t,
    \qquad
    y\left(\frac{\pi}{2}\right)=1,
    \qquad
    t>0.
\]
\end{problem}

\begin{solution}
Dividing by $t$ gives the standard form
\[
    y'+\frac{2}{t}y=\frac{2\sin t}{t}.
\]
Since $t>0$, an integrating factor is
\[
    \mu(t)=\exp\!\left(\int\frac{2}{t}\,\dd t\right)=t^2.
\]
Thus
\[
    (t^2y)'=2t\sin t.
\]
Integrating by parts gives
\[
    t^2y=-2t\cos t+2\sin t+c.
\]
Using $y(\pi/2)=1$ yields
\[
    \frac{\pi^2}{4}=2+c,
    \qquad
    c=\frac{\pi^2}{4}-2.
\]
Therefore
\[
    \boxed{
        y(t)=\frac{-2t\cos t+2\sin t+\frac{\pi^2}{4}-2}{t^2},
        \qquad t>0
    }.
\]
\end{solution}

\begin{problem}{2.2.5}
Solve the differential equation
\[
    \frac{\dd y}{\dd x}=\frac{3x-e^{-x}}{2y+e^y}.
\]
\end{problem}

\begin{solution}
Separating the variables gives
\[
    (2y+e^y)\,\dd y=(3x-e^{-x})\,\dd x.
\]
Integrating both sides,
\[
    y^2+e^y=\frac{3}{2}x^2+e^{-x}+c.
\]
Equivalently, the solution is given implicitly by
\[
    \boxed{2y^2-3x^2+2e^y-2e^{-x}=c},
    \qquad
    2y+e^y\neq0.
\]
\end{solution}

\begin{problem}{2.2.9}
Consider the initial value problem
\[
    y'=(1-2x)y^2,
    \qquad
    y(0)=-\frac{1}{12}.
\]
\begin{problemsubparts}
    \item Find the solution explicitly.
    \item Plot the graph of the solution.
    \item Determine (at least approximately) the interval in which the solution is defined.
\end{problemsubparts}
\end{problem}

\begin{solution}
\noindent\textup{(a)} Separating variables gives
\[
    y^{-2}\,\dd y=(1-2x)\,\dd x.
\]
Thus
\[
    -\frac{1}{y}=x-x^2+c.
\]
The initial condition gives $c=12$, so
\[
    \boxed{y(x)=\frac{1}{x^2-x-12}
    =\frac{1}{(x+3)(x-4)}}.
\]

\Needspace{23\baselineskip}
\medskip
\noindent\textup{(b)} The graph of the solution branch through $(0,-1/12)$ is shown below. The marked point is the initial point, and the dashed lines indicate the vertical asymptotes $x=-3$ and $x=4$.

\begin{center}
\begin{tikzpicture}
\begin{axis}[
    width=0.76\textwidth,
    height=0.42\textwidth,
    axis lines=middle,
    xmin=-3.5,
    xmax=4.5,
    ymin=-1.1,
    ymax=0.12,
    xlabel={$x$},
    ylabel={$y$},
    xtick={-3,0,4},
    ytick={-1,-0.5,0},
    tick label style={font=\small},
    label style={font=\small},
    clip=false
]
    \addplot[black,thick,domain=-2.93:3.93,samples=500] {1/(x^2-x-12)};
    \addplot[black!55,dashed] coordinates {(-3,-1.1) (-3,0.12)};
    \addplot[black!55,dashed] coordinates {(4,-1.1) (4,0.12)};
    \addplot[only marks,mark=*,mark size=1.8pt,black] coordinates {(0,-0.0833333333)};
\end{axis}
\end{tikzpicture}
\end{center}

\noindent\textup{(c)} The denominator vanishes at $x=-3$ and $x=4$. Therefore the largest open interval containing the initial point $x=0$ on which the solution is defined is
\[
    \boxed{(-3,4)}.
\]
\end{solution}

\begin{problem}{2.3.1}
Consider a tank used in certain hydrodynamic experiments. After one experiment the tank contains $150\,\mathrm{L}$ of a dye solution with a concentration of $1\,\mathrm{g/L}$. To prepare for the next experiment, the tank is rinsed with fresh water flowing in at a rate of $2\,\mathrm{L/min}$, while the well-stirred solution flows out at the same rate. Find the time that will elapse before the concentration of dye in the tank reaches $1\%$ of its original value.
\end{problem}

\begin{solution}
Let $Q(t)$ be the amount of dye in the tank, in grams, at time $t$ minutes. Initially,
\[
    Q(0)=150.
\]
Since fresh water contains no dye, the dye enters at rate $0$. The volume remains $150\,\mathrm{L}$, so the dye leaves at rate
\[
    \frac{Q(t)}{150}\cdot2=\frac{Q(t)}{75}\ \mathrm{g/min}.
\]
Hence
\[
    Q'=-\frac{1}{75}Q,
    \qquad
    Q(0)=150,
\]
and therefore
\[
    Q(t)=150e^{-t/75}.
\]
The concentration is $Q(t)/150=e^{-t/75}\,\mathrm{g/L}$. It reaches $1\%$ of its original value when
\[
    e^{-t/75}=0.01.
\]
Thus
\[
    \boxed{t=75\ln100\ \mathrm{min}\approx345.4\ \mathrm{min}}.
\]
\end{solution}

\begin{problem}{2.3.10}
Newton's law of cooling states that the temperature of an object changes at a rate proportional to the difference between its temperature and that of its surroundings. Suppose that the temperature of a cup of coffee obeys Newton's law of cooling. If the coffee has a temperature of $90^{\circ}\mathrm{C}$ when freshly poured, and $1$ minute later has cooled to $85^{\circ}\mathrm{C}$ in a room at $20^{\circ}\mathrm{C}$, determine when the coffee reaches a temperature of $65^{\circ}\mathrm{C}$.
\end{problem}

\begin{solution}
Let $T(t)$ be the temperature of the coffee, in degrees Celsius, at time $t$ minutes. Newton's law gives
\[
    T'=-k(T-20),
    \qquad
    k>0.
\]
Thus
\[
    T(t)=20+ce^{-kt}.
\]
Since $T(0)=90$, we have $c=70$, so
\[
    T(t)=20+70e^{-kt}.
\]
Using $T(1)=85$ gives
\[
    e^{-k}=\frac{65}{70}=\frac{13}{14}.
\]
When $T(t)=65$,
\[
    e^{-kt}=\frac{45}{70}=\frac{9}{14}.
\]
Therefore
\[
    \boxed{
        t=\frac{\ln(9/14)}{\ln(13/14)}\ \mathrm{min}
        \approx5.96\ \mathrm{min}
    }.
\]
\end{solution}

\end{document}
